\documentclass[11pt]{book}
\usepackage[paperwidth=145mm,paperheight=200mm,top=18mm,bottom=20mm,inner=20mm,outer=16mm]{geometry}
\usepackage{brumfiel}

% ---- local macros (hoist into book preamble) ----
% The book's run-in "Problem." paragraphs.
\newenvironment{IVproblem}{\par\medskip\noindent\emph{Problem.}\ }{\par\medskip}
% Centred display headings ("REVIEW OF CHAPTER 4", "ALGEBRA REVIEW").
\newcommand{\IVhead}[1]{%
  \par\bigskip\begin{center}\bfseries\large\MakeUppercase{#1}\end{center}%
  \par\nopagebreak\medskip}
% ---- end local macros ----

% figures04.tex -- Chapter 4 figures redrawn in TikZ.
%
% Label clearance, chapter-local (2026-08-17).  Every figure in this chapter
% is a point-on-a-line diagram whose labels sit directly over their point.
% The global `outer sep' was tightened to 2.2pt in the same pass that turned
% the point dots into fixed 3.2pt nodes -- and a 3.2pt dot reaches 1.6pt out
% of its point, so at 2.2pt the label box lands ON the dot it labels.
% check_labels bears this out exactly: the collision count for this chapter
% falls off a cliff at 3.2pt (39 -> 2) and does not improve above it.  3.2pt
% is therefore the tightest clearance that still clears the new dots, and is
% still tighter than the 3.4pt this book used before the dots existed.
% Coordinates are reconstructions; every congruence the printed figure
% asserts holds EXACTLY here (see tools/constraints/ch04.py).
% Self-contained: congruence ticks are drawn as explicit strokes, sized so
% that each prints at about 2.2pt half-length at the figure's own scale
% (Figs. 4-7/4-8/4-9, split out and enlarged, carry the book's own
% tick-to-segment ratio instead and so print nearer 3pt).
\tikzset{vlab/.append style={outer sep=3.2pt},
         slab/.append style={outer sep=3.2pt}}

% ------------------------------------------------------------- Fig 4-1
% Pictorial vignette (boy beside a marked spear stuck in hatched ground).
% Figurative line art, not a geometric diagram -> placeholder, per the
% plate rule. See ch04-UNCERTAIN.md.
\newcommand{\FIGIVONE}{\begin{bkfigure}[0.40\linewidth]
\includegraphics[width=0.40\linewidth]{fig-4-1}%
\figcap{4--1}\end{bkfigure}}

% -------------------------------------------------------- Figs 4-2, 4-3
% 4-2: III-1, the existence postulate.  AB congruent to A'B' (exact).
% 4-3: III-3, addition/subtraction.  AB=A'B' and BC=B'C' (exact: the two
%      lines carry identical parameter fractions over equal lengths).
\newcommand{\FIGIVTWOTHREE}{\begin{bkfigure}[\linewidth]
\begin{minipage}{0.47\linewidth}\centering
\begin{tikzpicture}[scale=0.86]
  \coordinate (A)  at (1.00000,0.42000);
  \coordinate (B)  at (4.20000,1.76000);
  \coordinate (Ap) at (1.55000,-0.53479);
  \coordinate (Bp) at (5.01280,-0.74609);
  \draw[fig] (0.12371,0.05306) -- (5.21465,2.18488);
  \draw[fig] (0.83242,-0.49100) -- (5.62546,-0.78347);
  \foreach \p in {A,B,Ap,Bp} {\dt{\p}}
  \node[vlab,above left]  at (A)  {$A$};
  \node[vlab,above]       at (B)  {$B$};
  \node[vlab,above left]  at (Ap) {$A'$};
  \node[vlab,above]       at (Bp) {$B'$};
  \node[vlab,right]       at (5.21465,2.18488) {$l$};
  \node[vlab,right]       at (5.62546,-0.78347) {$l'$};
\end{tikzpicture}
\figcap{4--2}
\end{minipage}\hfill
\begin{minipage}{0.47\linewidth}\centering
\begin{tikzpicture}[scale=0.86]
  \coordinate (A)  at (0.86972,0.09602);
  \coordinate (B)  at (2.70855,0.29902);
  \coordinate (C)  at (3.41923,0.37748);
  \coordinate (Ap) at (1.09193,-1.07678);
  \coordinate (Bp) at (2.93545,-0.92196);
  \coordinate (Cp) at (3.64794,-0.86212);
  % Left-hand stubs pulled in (2026-08-17): both lines ran 0.47 of AB to the
  % left of A / A', where the plate runs only about 0.22.  New start points
  % are exactly on the same lines (slope unchanged), so the incidence checks
  % in tools/constraints/ch04.py still hold -- their anchors move with them.
  \draw[fig] (0.46519,0.05136) -- (4.96981,0.54866);
  \draw[fig] (0.68636,-1.11084) -- (5.20247,-0.73157);
  \foreach \p in {A,B,C,Ap,Bp,Cp} {\dt{\p}}
  \node[vlab,above]      at (A)  {$A$};
  \node[vlab,above]      at (B)  {$B$};
  \node[vlab,above]      at (C)  {$C$};
  \node[vlab,above]      at (Ap) {$A'$};
  \node[vlab,above]      at (Bp) {$B'$};
  \node[vlab,above]      at (Cp) {$C'$};
  \node[vlab,right]      at (4.96981,0.54866)  {$l$};
  \node[vlab,right]      at (5.20247,-0.73157) {$l'$};
\end{tikzpicture}
\figcap{4--3}
\end{minipage}
\end{bkfigure}}

% ------------------------------------------------------------- Fig 4-4
% Exercise 4-2 #7.  AB=BC, CD=CE=ED (equilateral CDE), EF=EB.
% Construction: A,C,D collinear; E is the equilateral apex on CD below AD;
% B on ray EC beyond C with AB=BC (which forces ABC equilateral); F on ray
% ED beyond D with EF=EB.
\newcommand{\FIGIVFOUR}{\begin{bkfigure}[0.62\linewidth]
\begin{tikzpicture}[scale=0.689]
  \coordinate (A) at (0,0);
  \coordinate (C) at (3.60000,0);
  \coordinate (D) at (6.50000,0);
  \coordinate (E) at (5.05000,-2.51147);
  \coordinate (B) at (1.80000,3.11769);
  \coordinate (F) at (8.30000,3.11769);
  \draw[fig] (A) -- (D);
  \draw[fig] (A) -- (B) -- (E) -- (F);
  \node[vlab,below left]  at (A) {$A$};
  \node[vlab,above]       at (B) {$B$};
  \node[vlab,below left]  at (C) {$C$};
  \node[vlab,above left]  at (D) {$D$};
  \node[vlab,below]       at (E) {$E$};
  \node[vlab,above]       at (F) {$F$};
\end{tikzpicture}
\figcap{4--4}
\end{bkfigure}}

% ------------------------------------------------------------- Fig 4-5
% Exercise 4-2 #15/#16.  (a) two congruent copies: AB=A'B' (double tick),
% AC=A'C' (single tick on the arc).  (b) CA=BD.  (c) AC=BD.
\newcommand{\FIGIVFIVE}{\begin{bkfigure}[\linewidth]
\begin{tikzpicture}[scale=1.040]
  % ---- (a) left copy
  \begin{scope}[shift={(0.60,2.90)}]
    \coordinate (A) at (0,0); \coordinate (B) at (1.95000,0);
    \coordinate (C) at (2.85000,0);
    \draw[fig] (A) -- (C);
    \draw[fig] (A) .. controls +(0,-0.82667) and +(0,-0.82667) .. (C);
    \draw[tick] (0.90000,-0.09700) -- (0.90000,0.09700);
    \draw[tick] (1.05000,-0.09700) -- (1.05000,0.09700);
    \draw[tick] (1.42500,-0.71700) -- (1.42500,-0.52300);
    \foreach \p in {A,B,C} {\dt{\p}}
    \node[vlab,above] at (A) {$A$};
    \node[vlab,above] at (B) {$B$};
    \node[vlab,above] at (C) {$C$};
  \end{scope}
  % ---- (a) right copy
  \begin{scope}[shift={(5.20,2.90)}]
    \coordinate (Ap) at (0,0); \coordinate (Bp) at (1.95000,0);
    \coordinate (Cp) at (2.85000,0);
    \draw[fig] (Ap) -- (Cp);
    \draw[fig] (Ap) .. controls +(0,-0.82667) and +(0,-0.82667) .. (Cp);
    \draw[tick] (0.90000,-0.09700) -- (0.90000,0.09700);
    \draw[tick] (1.05000,-0.09700) -- (1.05000,0.09700);
    \draw[tick] (1.42500,-0.71700) -- (1.42500,-0.52300);
    \foreach \p in {Ap,Bp,Cp} {\dt{\p}}
    \node[vlab,above] at (Ap) {$A'$};
    \node[vlab,above] at (Bp) {$B'$};
    \node[vlab,above] at (Cp) {$C'$};
  \end{scope}
  \node[slab,left] at (0.20,2.90) {(a)};
  % ---- (b)
  \begin{scope}[shift={(0.60,1.45)}]
    \coordinate (Cb) at (1.51900,0); \coordinate (Ab) at (2.26900,0);
    \coordinate (Bb) at (4.67600,0); \coordinate (Db) at (5.42600,0);
    \draw[fig] (0,0) -- (7.00000,0);
    \draw[tick] (1.89400,-0.09700) -- (1.89400,0.09700);
    \draw[tick] (5.05100,-0.09700) -- (5.05100,0.09700);
    \foreach \p in {Cb,Ab,Bb,Db} {\dt{\p}}
    \node[vlab,below] at (Cb) {$C$};
    \node[vlab,below] at (Ab) {$A$};
    \node[vlab,below] at (Bb) {$B$};
    \node[vlab,below] at (Db) {$D$};
  \end{scope}
  \node[slab,left] at (0.20,1.45) {(b)};
  % ---- (c)
  \begin{scope}[shift={(0.60,0)}]
    \coordinate (Ac) at (0.75000,0); \coordinate (Bc) at (3.07000,0);
    \coordinate (Cc) at (4.07000,0); \coordinate (Dc) at (6.39000,0);
    \draw[fig] (0,0) -- (7.00000,0);
    \draw[fig] (Ac) .. controls +(0,0.96000) and +(0,0.96000) .. (Cc);
    \draw[fig] (Bc) .. controls +(0,0.96000) and +(0,0.96000) .. (Dc);
    \draw[tick] (2.41000,0.62300) -- (2.41000,0.81700);
    \draw[tick] (4.73000,0.62300) -- (4.73000,0.81700);
    \foreach \p in {Ac,Bc,Cc,Dc} {\dt{\p}}
    \node[vlab,below] at (Ac) {$A$};
    \node[vlab,below] at (Bc) {$B$};
    \node[vlab,below] at (Cc) {$C$};
    \node[vlab,below] at (Dc) {$D$};
  \end{scope}
  \node[slab,left] at (0.20,0) {(c)};
\end{tikzpicture}
\figcap{4--5}
\end{bkfigure}}

% ------------------------------------------------------------- Fig 4-6
% Definition 4-1.  AB congruent to A'B'' (exact); B' between A' and B''.
\newcommand{\FIGIVSIX}{\begin{bkfigure}[0.72\linewidth]
\begin{tikzpicture}[scale=1.170]
  \coordinate (A)   at (0.90000,1.32000);
  \coordinate (B)   at (4.55000,2.28000);
  \coordinate (Ap)  at (0.85000,0);
  \coordinate (Bp)  at (3.32023,0.09490);
  \coordinate (Bpp) at (4.62134,0.14489);
  \coordinate (Le)  at (5.47071,0.17752);
  \draw[fig] (A) -- (B);
  \draw[fig] (Ap) -- (Le);
  \foreach \p in {A,B,Ap,Bp,Bpp,Le} {\dt{\p}}
  \node[vlab,above left]  at (A)   {$A$};
  \node[vlab,above]       at (B)   {$B$};
  \node[vlab,above left]  at (Ap)  {$A'$};
  \node[vlab,above]       at (Bp)  {$B'$};
  \node[vlab,above]       at (Bpp) {$B''$};
  \node[vlab,above right] at (Le)  {$l'$};
\end{tikzpicture}
\figcap{4--6}
\end{bkfigure}}

% ------------------------------------------------ Figs 4-7, 4-8, 4-9
% 4-7 (Ex 4-3 #1): AB congruent to CB', with C-B'-D.
% 4-8 (Ex 4-3 #2): RS congruent to TS', with T-U-S'.
% 4-9 (Ex 4-3 #3): AB=DE (single tick), BC=EC' (double tick).
%
% SPLIT 2026-08-17: these three were one side-by-side block, printed after
% the whole of Exercise Group 4-3.  The book sets each one immediately after
% the exercise that cites it (4-7 after ex. 1 and 4-8 after ex. 2 in the left
% column of p.68, 4-9 above ex. 3 at the head of the right column), so each
% is now its own full-measure figure and ch04.tex places it inline.  Point
% coordinates are unchanged -- only the wrapper and `scale' moved -- so
% tools/constraints/ch04.py still describes them exactly.  The three scales
% are chosen so all three print about 4.4 cm wide, as the book draws them at
% one size; ticks ride the scale, keeping the book's tick-to-segment ratio.
\newcommand{\FIGIVSEVEN}{\begin{bkfigure}[0.56\linewidth]
\begin{tikzpicture}[scale=1.00]
  \coordinate (A)  at (1.20000,0.95000);
  \coordinate (B)  at (2.40000,0.95000);
  \coordinate (C)  at (0.92000,0);
  \coordinate (Bp) at (2.12000,0);
  \coordinate (D)  at (3.34000,0);
  \draw[fig] (0.30000,0.95000) -- (3.35000,0.95000);
  \draw[fig] (0,0) -- (4.40000,0);
  \draw[tick] (1.80000,0.84200) -- (1.80000,1.05800);
  \draw[tick] (1.52000,-0.10800) -- (1.52000,0.10800);
  \foreach \p in {A,B,C,Bp,D} {\dt{\p}}
  \node[vlab,above] at (A)  {$A$};
  \node[vlab,above] at (B)  {$B$};
  \node[vlab,below] at (C)  {$C$};
  \node[vlab,below] at (Bp) {$B'$};
  \node[vlab,below] at (D)  {$D$};
\end{tikzpicture}
\figcap{4--7}
\end{bkfigure}}

\newcommand{\FIGIVEIGHT}{\begin{bkfigure}[0.56\linewidth]
\begin{tikzpicture}[scale=1.20]
  \coordinate (R)  at (0.65000,1.30000);
  \coordinate (S)  at (3.00000,1.30000);
  \coordinate (T)  at (0.62000,0);
  \coordinate (U)  at (2.32000,0);
  \coordinate (Sp) at (2.97000,0);
  \draw[fig] (0,1.30000) -- (3.60000,1.30000);
  \draw[fig] (0,0) -- (3.65000,0);
  \draw[fig] (R) .. controls +(0,0.64000) and +(0,0.64000) .. (S);
  \draw[fig] (T) .. controls +(0,0.64000) and +(0,0.64000) .. (Sp);
  % tick half-length 0.090 here (not 0.108): this picture is drawn at
  % scale 1.20, so the shorter coordinate prints the same ink as 4-7/4-9.
  \draw[tick] (1.82500,1.69000) -- (1.82500,1.87000);
  \draw[tick] (1.79500,0.39000) -- (1.79500,0.57000);
  \foreach \p in {R,S,T,U,Sp} {\dt{\p}}
  \node[vlab,above left]  at (R)  {$R$};
  \node[vlab,above right] at (S)  {$S$};
  \node[vlab,below] at (T)  {$T$};
  \node[vlab,below] at (U)  {$U$};
  \node[vlab,below] at (Sp) {$S'$};
\end{tikzpicture}
\figcap{4--8}
\end{bkfigure}}

\newcommand{\FIGIVNINE}{\begin{bkfigure}[0.56\linewidth]
\begin{tikzpicture}[scale=1.02]
  \coordinate (A)  at (0.75000,0.95000);
  \coordinate (B)  at (2.60000,0.95000);
  \coordinate (C)  at (3.40000,0.95000);
  \coordinate (D)  at (0.75000,0);
  \coordinate (E)  at (2.60000,0);
  \coordinate (Cp) at (3.40000,0);
  \draw[fig] (0,0.95000) -- (4.30000,0.95000);
  \draw[fig] (0,0) -- (4.30000,0);
  \draw[tick] (1.67500,0.84200) -- (1.67500,1.05800);
  \draw[tick] (2.91700,0.84200) -- (2.91700,1.05800);
  \draw[tick] (3.08300,0.84200) -- (3.08300,1.05800);
  \draw[tick] (1.67500,-0.10800) -- (1.67500,0.10800);
  \draw[tick] (2.91700,-0.10800) -- (2.91700,0.10800);
  \draw[tick] (3.08300,-0.10800) -- (3.08300,0.10800);
  \foreach \p in {A,B,C,D,E,Cp} {\dt{\p}}
  \node[vlab,above] at (A)  {$A$};
  \node[vlab,above] at (B)  {$B$};
  \node[vlab,above] at (C)  {$C$};
  \node[vlab,below] at (D)  {$D$};
  \node[vlab,below] at (E)  {$E$};
  \node[vlab,below] at (Cp) {$C'$};
\end{tikzpicture}
\figcap{4--9}
\end{bkfigure}}

% ------------------------------------------------------ Figs 4-10, 4-11
% 4-10 (Thm 4-2): AB=A'C' (double tick on the arcs), AC=A'B' (single tick).
% 4-11: AC=A'B' (single tick), CX=B'C' (double tick).
\newcommand{\FIGIVTENELEVEN}{\begin{bkfigure}[\linewidth]
\begin{minipage}{0.47\linewidth}\centering
\begin{tikzpicture}[scale=0.901]
  \coordinate (A)  at (0.30000,1.15000);
  \coordinate (C)  at (2.75000,1.15000);
  \coordinate (B)  at (3.45000,1.15000);
  \coordinate (Ap) at (0.30000,0);
  \coordinate (Bp) at (2.75000,0);
  \coordinate (Cp) at (3.45000,0);
  % Right-hand stubs trimmed 4.30 -> 3.80 (2026-08-17).  Measured off the
  % plate, each line runs about 0.10 of AB past B / C'; ours ran 0.27 past,
  % so both lines were over-drawn to the right.  The left stubs (0.30, also
  % 0.10 of AB) already matched.
  \draw[fig] (0,1.15000) -- (3.80000,1.15000);
  \draw[fig] (0,0) -- (3.80000,0);
  \draw[fig] (A)  .. controls +(0,0.76000) and +(0,0.76000) .. (B);
  \draw[fig] (Ap) .. controls +(0,-0.76000) and +(0,-0.76000) .. (Cp);
  \draw[tick] (1.80000,1.62300) -- (1.80000,1.81700);
  \draw[tick] (1.95000,1.62300) -- (1.95000,1.81700);
  \draw[tick] (1.80000,-0.66700) -- (1.80000,-0.47300);
  \draw[tick] (1.95000,-0.66700) -- (1.95000,-0.47300);
  \draw[tick] (1.52500,1.05300) -- (1.52500,1.24700);
  \draw[tick] (1.52500,-0.09700) -- (1.52500,0.09700);
  % The two correspondence arrows.  In the book they bow AWAY from each
  % other, so that "then (" and ") if" read as a facing pair of parentheses:
  % the B'->C arrow bows left (west), the B->C' arrow bows right (east).
  % Both were drawn with `bend right' here, which mirrored each one and made
  % the pair read ")(" instead.  For an upward path `bend left' is the
  % westward bow; for a downward path it is the eastward one.
  \draw[fig,->,line width=0.55pt] (2.72000,0.14000)
        to[bend left=28] (2.72000,1.01000);
  \draw[fig,->,line width=0.55pt] (3.45000,1.01000)
        to[bend left=28] (3.45000,0.14000);
  \foreach \p in {A,C,B,Ap,Bp,Cp} {\dt{\p}}
  \node[vlab,above left]  at (A)  {$A$};
  \node[vlab,above]       at (C)  {$C$};
  % B and C' are the far ends of the two arcs, so the arc arrives into the
  % label's own corner.  The plate sets both in the free quadrant on the
  % outside of the arc -- B above-RIGHT of its dot, C' below-RIGHT -- which
  % is also what keeps them clear of the arc ends here.
  \node[vlab,above right] at (B)  {$B$};
  \node[vlab,below left]  at (Ap) {$A'$};
  \node[vlab,below]       at (Bp) {$B'$};
  \node[vlab,below right] at (Cp) {$C'$};
  \node[slab,left]  at (2.42000,0.57500) {then};
  \node[slab,right] at (3.78000,0.57500) {if};
\end{tikzpicture}
\figcap{4--10}
\end{minipage}\hfill
\begin{minipage}{0.47\linewidth}\centering
\begin{tikzpicture}[scale=0.901]
  \coordinate (A)  at (0.60000,0.95000);
  \coordinate (C)  at (2.55000,0.95000);
  \coordinate (X)  at (3.35000,0.95000);
  \coordinate (Ap) at (0.60000,0);
  \coordinate (Bp) at (2.55000,0);
  \coordinate (Cp) at (3.35000,0);
  \draw[fig] (0,0.95000) -- (4.20000,0.95000);
  \draw[fig] (0,0) -- (4.20000,0);
  \draw[tick] (1.57500,0.85300) -- (1.57500,1.04700);
  \draw[tick] (2.87500,0.85300) -- (2.87500,1.04700);
  \draw[tick] (3.02500,0.85300) -- (3.02500,1.04700);
  \draw[tick] (1.57500,-0.09700) -- (1.57500,0.09700);
  \draw[tick] (2.87500,-0.09700) -- (2.87500,0.09700);
  \draw[tick] (3.02500,-0.09700) -- (3.02500,0.09700);
  \foreach \p in {A,C,X,Ap,Bp,Cp} {\dt{\p}}
  \node[vlab,above] at (A)  {$A$};
  \node[vlab,above] at (C)  {$C$};
  \node[vlab,above] at (X)  {$X$};
  \node[vlab,below] at (Ap) {$A'$};
  \node[vlab,below] at (Bp) {$B'$};
  \node[vlab,below] at (Cp) {$C'$};
\end{tikzpicture}
\figcap{4--11}
\end{minipage}
\end{bkfigure}}

% ------------------------------------------------------ Figs 4-12, 4-13
% 4-12 (Thm 4-3): AB=A'B' (double tick), AC=A'C' (single tick on the arcs).
% 4-13 (Thm 4-4): AD'=CD=EF.
\newcommand{\FIGIVTWELVETHIRTEEN}{\begin{bkfigure}[\linewidth]
\begin{minipage}{0.47\linewidth}\centering
\begin{tikzpicture}[scale=0.833]
  \coordinate (A)  at (0.30000,1.25000);
  \coordinate (B)  at (2.75000,1.25000);
  \coordinate (C)  at (3.50000,1.25000);
  \coordinate (Ap) at (0.30000,0);
  \coordinate (Bp) at (2.75000,0);
  \coordinate (Cp) at (3.50000,0);
  % Right-hand stubs trimmed 4.10 -> 3.90: the plate runs about 0.12 of AC
  % past C / C', ours ran 0.19 past.
  \draw[fig] (0,1.25000) -- (3.90000,1.25000);
  \draw[fig] (0,0) -- (3.90000,0);
  \draw[fig] (A)  .. controls +(0,-0.60000) and +(0,-0.60000) .. (C);
  \draw[fig] (Ap) .. controls +(0,-0.60000) and +(0,-0.60000) .. (Cp);
  \draw[tick] (1.45000,1.15300) -- (1.45000,1.34700);
  \draw[tick] (1.60000,1.15300) -- (1.60000,1.34700);
  \draw[tick] (1.45000,-0.09700) -- (1.45000,0.09700);
  \draw[tick] (1.60000,-0.09700) -- (1.60000,0.09700);
  \draw[tick] (1.90000,0.70300) -- (1.90000,0.89700);
  \draw[tick] (1.90000,-0.54700) -- (1.90000,-0.35300);
  % The B'->B correspondence arrow.  The plate draws a full arrow from just
  % above B' up to just under B, crossing the upper arc on the way; ours was
  % a short stub hanging off B and pointing the wrong way (down into empty
  % space), so the figure said nothing about which point matched which.
  % Upward path + `bend right' = the plate's slight eastward bow.
  \draw[fig,->,line width=0.55pt] (2.79000,0.15000)
        to[bend right=14] (2.75000,1.11000);
  \foreach \p in {A,B,C,Ap,Bp,Cp} {\dt{\p}}
  \node[vlab,above] at (A)  {$A$};
  \node[vlab,above] at (B)  {$B$};
  \node[vlab,above] at (C)  {$C$};
  \node[vlab,below left] at (Ap) {$A'$};
  \node[vlab,below]      at (Bp) {$B'$};
  \node[vlab,below]      at (Cp) {$C'$};
\end{tikzpicture}
\figcap{4--12}
\end{minipage}\hfill
\begin{minipage}{0.47\linewidth}\centering
\begin{tikzpicture}[scale=0.833]
  \coordinate (A)  at (0,1.45000);
  \coordinate (Dp) at (3.34012,1.70701);
  \coordinate (B)  at (5.63333,1.88347);
  \coordinate (C)  at (0,0.72000);
  \coordinate (D)  at (3.33824,1.00033);
  \coordinate (E)  at (0,0);
  \coordinate (F)  at (3.34776,0.12274);
  \draw[fig] (A) -- (B);
  \draw[fig] (C) -- (D);
  \draw[fig] (E) -- (F);
  \foreach \p in {A,Dp,B,C,D,E,F} {\dt{\p}}
  \node[vlab,above left] at (A)  {$A$};
  \node[vlab,above]      at (Dp) {$D'$};
  \node[vlab,above]      at (B)  {$B$};
  \node[vlab,above left] at (C)  {$C$};
  \node[vlab,above]      at (D)  {$D$};
  \node[vlab,above left] at (E)  {$E$};
  \node[vlab,above]      at (F)  {$F$};
\end{tikzpicture}
\figcap{4--13}
\end{minipage}
\end{bkfigure}}


\begin{document}
\setcounter{chapter}{4}

\chapter*{\raggedright\Huge CONGRUENCE OF SEGMENTS}
\addcontentsline{toc}{chapter}{4.\ Congruence of Segments}
\markboth{CONGRUENCE OF SEGMENTS}{}
\vspace{-1em}\noindent\rule{\linewidth}{0.6pt}\vspace{1em}
\figurekey

%========================================================= 4-1
\section*{4--1\quad INTRODUCTION}
\addcontentsline{toc}{section}{4--1\quad Introduction}

In this chapter we shall be concerned with comparing the ``sizes'' of segments.
All of us have learned to understand the terms ``distance'' and ``length'' as
used in everyday life. For example, we may say that it is 12 miles to the
nearest town, or that a fence is 153 feet long. In each case we give a number to
indicate a distance or length. To give a number is a most convenient means of
describing the distance between two points or, in other words, of telling how far
apart they are.

But even before men learned to count, they must have wished at times to tell how
far apart points were, or to decide if two objects had the same length. For
example, a hunter visiting a friend in a neighboring village might have said to
his friend, ``My son is as tall as this mark on my spear.'' Then the men could
see, by comparing with the mark, whether the host's son was as tall as the
visitor's son. In this example the marked spear served as a kind of
``yardstick'' for comparing the height of the hunter's son with the height of
others. Even in very primitive societies men must have had a means for comparing
the sizes of objects.

\FIGIVONE

\begin{IVproblem}
Describe other examples where objects are compared in size without the use of
measurement, that is, a number. Some of your examples may use a marked stick (or
``yardstick'') and some not.
\end{IVproblem}

We are so used to speaking of distance or of length in terms of numbers that it
is hard for us to find other words to describe how far apart two points are.
Although we usually associate numbers with the words ``distance'' and
``length,'' what we study in this chapter is so simple that it would be
understandable to a person who knew nothing of numbers at all. We will be talking
about the ``size'' of a segment and whether two segments are ``the same size.''
As we have seen with the example of the hunter, a practical method for comparing
segments is to use a marked stick and to carry it from one place to another. We
shall see in the Section 4--2 how in geometry we can avoid carrying such sticks.

\exercisegroup{4--1}

\begin{multicols}{2}
\begin{exlist}
\item Your foot is a natural ``yardstick'' you carry with you. Describe several
      examples of using your foot for comparing lengths.
\item Think of several ways to use your hands as ``yardsticks.''
\end{exlist}
\end{multicols}

%========================================================= 4-2
\section*{4--2\quad CONGRUENCE}
\addcontentsline{toc}{section}{4--2\quad Congruence}

In daily life we are often interested in size, and we have ways of testing
whether or not two things are the same size. But our geometry so far has in it
\emph{no notion of size at all}. Because our geometry deals with imaginary
concepts (points and lines), to talk about size it is necessary to assume
(postulate) properties that we desire. We do not wish to talk about ``carrying a
segment around in the plane,'' because this involves the idea of movement, which
is not a part of geometry proper. Consequently \emph{we will just suppose that by
some means (it does not matter how) we can tell whether or not two segments are
the same size.} In other words, we will assume that we can pick out all the
segments of the same size.

We have been using the words ``the same size.'' We now use another word which is
to have exactly the same intuitive significance. The word is \emph{congruent},
and whenever you see it you can, if you like, say to yourself, ``the same size.''
Congruence refers to the \emph{relation} of being the same size. The word comes
from a Latin word meaning ``agree'' or ``conform.''

Euclid used the word ``equal'' instead of congruent. However, \emph{in this book
we will use the word ``equal'' (and the sign ``$=$'') to mean ``identically the
same.''} Thus, if we write $A = B$ we mean that $A$ and $B$ are one and the same
point. Or, if $x$ and $y$ are numbers, $x = y$ means that they are one and the
same number. And in dealing with segments, if we write $AB = CD$, we mean that
$A = C$ and $B = D$ or $A = D$ and $B = C$, that is, $AB$ and $CD$ are one and
the same segment.

%========================================================= 4-3
\section*{4--3\quad POSTULATE GROUP III. LINEAR CONGRUENCE}
\addcontentsline{toc}{section}{4--3\quad Postulate Group III. Linear Congruence}

We assume that there is a relation between segments called \emph{congruence}.
This relation is to have the following properties.

III--1 \emph{The existence postulate for segments.} Given points $A$ and $B$ on a
line $l$ and a point $A'$ on a line $l'$ (Fig.~4--2), then on each side of $A'$
on $l'$ there is exactly one point $B'$ such that $AB$ is congruent to $A'B'$. To
indicate that $AB$ is congruent to $A'B'$, we write ``$AB \cong A'B'$.''

\FIGIVTWOTHREE

In congruence the order of the points does not matter, that is, if
$AB \cong A'B'$, then also $AB \cong B'A'$, $BA \cong A'B'$, and
$BA \cong B'A'$.

III--2 \emph{The three-part segment postulate.} Congruence also satisfies the
following:

\begin{enumerate}[label=(\alph*), leftmargin=2.2em, itemsep=2pt, topsep=3pt]
\item $AB \cong AB$. (A segment is congruent to itself; that is, congruence of
      segments is reflexive.)
\item If $AB \cong A'B'$, then $A'B' \cong AB$. (Congruence of segments is
      symmetric.)
\item If $AB \cong A'B'$ and $A'B' \cong A''B''$, then $AB \cong A''B''$.
      (Congruence of segments is transitive.)
\end{enumerate}

\begin{IVproblem}
Prove that if $AB \cong CD$ and $EF \cong CD$, then $AB \cong EF$.
[\emph{Hint:} Use parts (b) and (c) of III--2.]
\end{IVproblem}

III--3 \emph{The addition and subtraction of segments postulate.} Suppose that
$B$ is between $A$ and $C$ on a line $l$, and that $B'$ is between $A'$ and $C'$
on a line $l'$ (Fig.~4--3).

\begin{enumerate}[label=(\alph*), leftmargin=2.2em, itemsep=1pt, topsep=3pt]
\item If $AB \cong A'B'$ and $BC \cong B'C'$, then $AC \cong A'C'$.
\item If $AC \cong A'C'$ and $BC \cong B'C'$, then $AB \cong A'B'$.
\end{enumerate}

These postulates are somewhat like our earlier ones. They seem so obvious that
you might feel they do not need to be stated. As before, however, you must
remember that points and lines are inventions of the mind, and so we must state
clearly what properties they are to have. Euclid did not state them as clearly as
we have. The postulates given in this book are easier to use than those of Euclid
because they tell you exactly what you can do. They are similar to ones given by
David Hilbert (see Section 1--2).

In less exact language, the postulates say the following things. The first says
that on the line $l'$ there are two points $B'$ and $B''$, one on each side of
$A'$, such that $AB$ is the same size as both $A'B'$ and $A'B''$. The second one
says three things: that a segment is the same size as itself; that if $AB$ is the
same size as $A'B'$, then $A'B'$ is the same size as $AB$; and that if two
segments are the same size as a third segment, then the two segments are the same
size. The third postulate tells how to combine segments by what we call
\emph{addition and subtraction of segments}.

\exercisegroup{4--2}

\begin{multicols}{2}
\begin{exlist}
\item What does it mean to say $RS = TU$?
\item If $RS = TU$, is $RS \cong TU$? Justify your answer.
\item Is it true for every pair of segments $RS$ and $TU$ that if
      $RS \cong TU$ then $RS = TU$?
\item Two segments $AB$ and $CD$ are not congruent to each other. (We write this
      as $AB \ncong CD$.) Prove that $AB$ and $CD$ are not equal to each other
      ($AB \neq CD$).
\item If $ABC$, what postulates must you use to prove that $AB \ncong AC$?
\item Is the following statement true or false? If $AB$ is any segment and $A'$
      is any point on a line $l'$, then there is one and only one point $B'$ on
      $l'$ such that $AB \cong A'B'$.
\item In Fig.~4--4 the following congruences hold: $AB \cong BC$,
      $CD \cong CE$, $CE \cong ED$, $EF \cong EB$. List all other congruences
      that you recognize.
\item The sign $=$ was defined in this section. State the definition in your own
      words.
\item State the three postulates of linear congruence in your own words, using
      the phrase ``the same size'' instead of ``congruent.'' Then go back and
      read again the paragraph preceding this group of exercises.
\item Show (using postulates) that if $A$, $B$, $C$, $D$ are points, in this
      order, on a line $l$, and if $A'$, $B'$, $C'$, $D'$ are points, in this
      order, on a line $l'$, and if $AB \cong A'B'$, $BC \cong B'C'$,
      $CD \cong C'D'$, then $AD \cong A'D'$. [\emph{Hint:} Draw a figure. You
      will have to use part (a) of Postulate III--3 twice.]
\item Assuming that the points have the same order as in Exercise 10, prove that
      $(AD \cong A'D'$ and $CD \cong C'D'$ and
      $BC \cong B'C') \rightarrow AB \cong A'B'$.
\item Suppose $E$, $F$, $G$ and $E'$, $F'$, $G'$ are \emph{any} points on lines
      $l$ and $l'$, respectively, such that $EF \cong E'F'$ and
      $FG \cong F'G'$. Show, by drawing a figure, that it does not need to be
      true that $EG \cong E'G'$. What do we need to know about the position of
      the points? Does III--3 take care of this?
\item[\stex 13.] Part (b) of Postulate III--3 can be \emph{proved} from the
      other postulates. We have included this as a postulate to make things
      easier. Try to prove (b). You will have to use (a) and Postulate III--1.
\item[\stex 14.] Try to reword III--3 in terms of length.
\item[15.] Each part of Fig.~4--5 indicates certain congruence relations. The
      marks show which segments are congruent. What other relations can be
      deduced? Point out how the postulates are used.
\item[\stex 16.] Show how to relabel the points in parts (b) and (c) of
      Fig.~4--5 as $A$, $B$, $C$ and $A'$, $B'$, $C'$ so that the application of
      Postulate III--3 is clearly shown.
\end{exlist}
\end{multicols}

\FIGIVFOUR

\FIGIVFIVE

%========================================================= 4-4
\section*{4--4\quad COMPARING SEGMENTS}
\addcontentsline{toc}{section}{4--4\quad Comparing Segments}

The postulates in Group III tell us about segments that are the same size or
congruent to each other. We can now \emph{define}, in terms of congruence and
betweenness, the ideas of larger and smaller segments. The definition will simply
describe what the hunter did when he compared the boy with the mark on the spear.

We need a symbol to stand for ``is greater than'' and another symbol to stand for
``is less than.'' These symbols are: $>$ and $<$. Whenever you see $>$, read
``is greater than,'' and whenever you see $<$, read ``is less than.''

\begin{definition}
Let $AB$ be any segment and $A'B'$ any other segment on line $l'$ (Fig.~4--6). On
the same side of $A'$ as $B'$, let $B''$ be the point such that
$AB \cong A'B''$. If $A'B'B''$ (that is, if $B'$ is between $A'$ and $B''$), we
say that $AB > A'B'$. If $A'B''B'$ (that is, $B''$ is between $A'$ and $B'$), we
say that $AB < A'B'$.
\end{definition}

\FIGIVSIX

\begin{remark}
Because of III--1 there is only one point $B''$ on the same side of $A'$ as the
point $B'$ and such that $AB \cong A'B''$. The three points $A'$, $B'$, and
$B''$ are therefore in some definite order (unless $B' = B''$, in which case
$AB \cong A'B'$). We cannot have $A'$ between $B'$ and $B''$ because $B'$ and
$B''$ are on the same side of $A'$. Therefore if $AB \ncong A'B'$, we must have
either $A'B'B''$ or $A'B''B'$. In other words, either $AB > A'B'$ or
$AB < A'B'$. This remark proves the following theorem.
\end{remark}

\begin{theorem}
If $AB$ and $A'B'$ are any two segments, then either $AB > A'B'$, or
$AB \cong A'B'$, or $AB < A'B'$, and exactly one of these is true.
\end{theorem}

\exercisegroup{4--3}

\begin{multicols}{2}
\begin{exlist}
\item Does Fig.~4--7 indicate that $AB < CD$ or that $CD > AB$? Read
      Definition 4--1 again.
\item What relationships are indicated by Fig.~4--8? Use Definition 4--1.
\item What relationships are indicated by Fig.~4--9? Use Definition 4--1.
\item If we know that $AB \cong CD$ and $ABX$, what can we conclude? Use
      Definition 4--1.
\item It is known that $RS \cong UT$ and $RVS$. What conclusion can be drawn?
      Use Definition 4--1.
\item Does Definition 4--1 say that if $AB > CD$ then $CD < AB$? Does looking at
      a figure prove that this is so?
\end{exlist}
\end{multicols}

\FIGIVSEVENEIGHTNINE

The exercises above are meant to prepare you for the main theorem of this
section, namely: if one segment is greater than a second, then the second is less
than the first. Of course, we know that if we draw two pencil marks so that the
first is longer than the second, then the second is shorter than the first. But
we have not yet defined \emph{length of a segment}. As a matter of fact we shall
need to introduce a new postulate into our geometry before length can be defined.
Since our definitions of ``greater than'' and ``less than'' have been formulated
in terms of betweenness, the truth of any statement about these relations
(``greater than'' and ``less than'') must be established by using properties of
betweenness.

% Theorem 4-2 carries a dagger and a marker-less footnote in the book, so it is
% set by hand rather than through the theorem environment; the counter is still
% stepped so that Theorem 4-3 numbers correctly.
\refstepcounter{theorem}%
\par\medskip\noindent
\textbf{$\dagger$Theorem \thetheorem.}\ \emph{If $AB > A'B'$ then
$A'B' < AB$, and conversely, if $A'B' < AB$ then $AB > A'B'$.}%
\begingroup\let\thefootnote\relax
\footnotetext{$\dagger$ Although this theorem is intuitively obvious, its proof
is quite difficult. If the teacher wishes, it may be treated as a postulate.}%
\endgroup
\par\medskip

\noindent\emph{Proof.} Figure 4--10 illustrates what must be proved in order to
establish the first part of the theorem.

\begin{enumerate}[label=(\arabic*), leftmargin=2.2em, itemsep=1pt, topsep=3pt]
\item $AB > A'B'$ means $AB \cong A'C'$ and $A'B'C'$.
\item $A'B' < AB$ means $A'B' \cong AC$ and $ACB$.
\end{enumerate}

\FIGIVTENELEVEN

To show that $AB > A'B' \rightarrow A'B' < AB$, we must prove that if $C'$ falls
``beyond'' $B'$, then $C$ falls between $A$ and $B$.

From (1) above we have $AB \cong A'C'$ and $A'B'C'$. We know also that on a
given side of $A$ there is a point $C$ such that $AC \cong A'B'$. To establish
that we do indeed have $ACB$, consider the point $X$ in Fig.~4--11 such that
$ACX$ and $CX \cong B'C'$. Now, by addition of segments, $AX \cong A'C'$. But
$AB \cong A'C'$, so $X = B$. (Why?) And since $ACX$, we have $ACB$; and from (2)
above we have $A'B' < AB$.

The proof of the converse is similar to the above proof and is left as an
exercise.

From now on we shall freely interchange the ``greater than'' and ``less than''
terminology since the statements $AB > CD$ and $CD < AB$ are equivalent to each
other.

\begin{theorem}
If $AC \cong A'C'$ and $A'B'C'$ and $B$ is the point on $C$'s side of $A$ such
that $AB \cong A'B'$, then $ABC$ (Fig.~4--12).
\end{theorem}

\noindent\emph{Proof.} Since $AC > A'B'$, $A'B' < AC$, and so by Definition 4--1
we have $B$ between $A$ and $C$.

\medskip
This theorem is proved at this time because we shall need it in the next chapter
when we discuss the length of a segment. We can then show that congruent segments
have the same length. It is also a very convenient tool to use in many subsequent
proofs.

\FIGIVTWELVETHIRTEEN

\begin{definition}
If we wish to say ``$AB \cong A'B'$ or $AB > A'B'$,'' we write
``$AB \geqq A'B'$'' and read the symbol $\geqq$ as ``is greater than or
congruent to.'' Similarly, if we wish to say ``$AB \cong A'B'$ or
$AB < A'B'$,'' we write ``$AB \leqq A'B'$'' and read the symbol $\leqq$ as ``is
less than or congruent to.''
\end{definition}

\begin{theorem}
If $AB > CD$ and $CD \cong EF$, then $AB > EF$.
\end{theorem}

\noindent\emph{Proof.} Since $CD < AB$ (why?) there is a point $D'$
(Fig.~4--13) on $AB$ such that $CD \cong AD'$. But then, by III--2,
$EF \cong AD'$. Hence, by Definition 4--1, $EF < AB$ and $AB > EF$.

\exercisegroup{4--4}

\begin{multicols}{2}
\begin{exlist}
\item Using the definition of $>$, explain what must be done in order to prove
      that $AB > CD$. Draw a figure to illustrate this.
\item What must one show in order to prove $RS < TU$? Draw a figure to
      illustrate this.
\item What basic undefined relations are used to define the phrases ``is greater
      than'' and ``is less than''?
\item Is it correct to say that $AB \geqq CD \rightarrow AB \ncong CD$? Is the
      converse of this implication true?
\item Is it correct to say that $AB > CD \rightarrow AB \ncong CD$? State both
      the converse and contrapositive of this statement. Which one of these
      statements is true?
\item What conclusion can you draw if you know that $A$, $B$, $C$, and $D$ are
      points on line $l$ and that $AB > CD$ and $A = C$?
\item If $AB$ is neither equal to nor greater than $CD$, can you conclude
      $AB < CD$?
\item If $AB$ is neither congruent to nor less than $CD$, can you conclude
      $AB > CD$?
\item[\stex 9.] Prove that if $AB < EF$ and $EF \cong CD$, then $AB < CD$.
      State the converse and contrapositive of this theorem.
\item[\stex 10.] Prove that if $AB > CD$ and $CD > EF$, then $AB > EF$.
\item[\stex 11.] Prove that if $AB < CD$ and $CD < EF$, then $AB < EF$.
\item[\stex 12.] Prove that if $B$ is between $A$ and $C$, and $B'$ is between
      $A'$ and $C'$, and $AB \cong A'B'$, and $BC > B'C'$, then $AC > A'C'$.
\item[13.] If $AB > CD$ and $CD < EF$, can you prove that $AB > EF$?
\item[\stex 14.] Prove that if $ABC$ and $A'C'B'$ and $AC \cong A'C'$, then
      $AB < A'B'$.
\end{exlist}
\end{multicols}

%========================================================= Review
\IVhead{Review of Chapter 4}
\addcontentsline{toc}{section}{Review of Chapter 4}
\markboth{REVIEW OF CHAPTER 4}{}

Several important ideas have been developed in this chapter. You should realize
that \emph{congruence} is really an undefined relation. Of course, you are so
accustomed to comparing the sizes of things by comparing numbers that it is
probably quite difficult to see clearly that numbers are not actually needed. We
know that a $6'8''$ basketball center is taller than a
$6'2''$ center. We know that a 240-lb football tackle is
larger than a 180-lb tackle. But such comparisons as these can be made
\emph{without using numbers at all}. Surely men were able to recognize
differences in sizes long before they had a number system to measure them.
Historically, the Greeks developed their geometry without the use of number
concepts.

In the next chapter, when we define what is meant by the length of a segment, it
will naturally turn out that congruent segments have the same length. That is, if
$AB \cong CD$, and $x$ and $y$ are numbers which are the lengths of these
segments, then $x = y$.

Again, the postulates of this chapter are so obvious to our intuition that you
might wonder why we bother to make them. All we can say is that postulates should
be obvious. The amazing thing about mathematics is that out of such simple
assumptions as we are making in our study of geometry one can construct theorems
of marvelous complexity which provide the foundation for all work in science.

The proofs in this chapter are quite difficult. There is no need to memorize the
results of the theorems since they merely show that our geometry agrees with the
figures that we draw.

In working with Definition 4--1, you should have learned an important
mathematical lesson, namely, that definitions must be clearly understood. It
would be very easy to read Definition 4--1 carelessly and to assume that if
$AB > CD$ then $CD < AB$. This is not stated in the definition, but of course we
were able to use the definition, our postulates, and the rules of logic to prove
that it is indeed true.

The exercises below will help you review the concepts of this chapter.

\medskip
Label the following statements true or false.

\begin{exlist}
\item If $AB \cong CD$, then $A = C$ and $B = D$.
\item If $A = C$ and $B = D$, then $AB \cong CD$.
\item If $AB \cong CD$ or $AB < CD$, then $AB \leqq CD$.
\item If $A$, $B$, and $C$ are points on a line $l$ and $A'$, $B'$, and $C'$ are
      points on a line $l'$, and $AB \cong A'B'$, and $BC \cong B'C'$, then
      $AC \cong A'C'$.
\item If $AB$ is any segment and $A'$ is any point on a line $l'$, then there are
      exactly two points $B'$ and $B''$ on $l'$ such that $A'B'$ and $A'B''$ are
      congruent to $AB$.
\item If $AB$ is not less than $CD$, then it must be greater than $CD$.
\item If $RS$ is neither greater than nor less than $TU$, then $RS \cong TU$.
\item If $A$, $B$, $C$, and $D$ are points in this order on a line $l$ and
      $AC \cong BD$, then $AB \cong CD$.
\item If $A$, $B$, $C$, and $D$ are points in this order on a line $l$ and
      $AC > BD$, then $AB < CD$.
\item If $A$, $B$, $C$, and $D$ are points in this order on a line $l$ and
      $AC < BD$, then $AB \leqq CD$.
\end{exlist}

\medskip
This chapter has included three postulates, two definitions, and four theorems.
The following statements are all true. Designate a postulate or definition or
theorem which could be used in justifying each statement.

\begin{exlist}[start=11]
\item For every line segment $AB$, it is true that $AB \cong AB$.
\item If $AB \cong CD$ and $ABX$, then $AX > CD$.
\item If $AB$ is not less than $CD$, then $AB \geqq CD$.
\item If $AB \cong CD$ and $CD \cong EF$ and $XY > AB$, then $XY > EF$.
\item If $RS \cong TU$ and $XY \cong TU$, then $RS \cong XY$.
\item If $ABC$ and $BCD$, then $AC < AD$.
\item If $AB \leqq CD$, then there is no point $E$ such that $CDE$ and
      $AB \cong CE$.
\item If in triangle $ABC$ we have side $AC$ greater than side $AB$, then there
      is a point $E$ on side $AC$ such that $AE \cong AB$.
\item If $ABC$ and $AB \cong BC$ and $XYZ$ and $XY \cong YZ \cong AB$, then
      $AC \cong XZ$.
\item If $A = X$ and $B = Y$, then $AY \cong BX$.
\end{exlist}

%========================================================= Algebra review
\IVhead{Algebra Review}
\addcontentsline{toc}{section}{Algebra Review}

In algebra the relations of ``greater than'' and ``less than'' are defined for
pairs of numbers. We say that $5 > 3$ because $5 = 3 + 2$, and 2 is a positive
number. More generally we make the following definitions:

\begin{enumerate}[label=(\arabic*), leftmargin=2.2em, itemsep=2pt, topsep=3pt]
\item If $a$, $b$, and $c$ are numbers and $a = b + c$ and $c$ is positive, then
      we say that $a$ is greater than $b$ and write $a > b$.
\item If $a$, $b$, and $c$ are numbers and $a = b - c$ and $c$ is a positive
      number, then we say that $a$ is less than $b$ and write $a < b$.
\end{enumerate}

For example, since $-\frac{3}{2} = -\frac{5}{2} + \frac{2}{2}$, we say
$-\frac{3}{2} > -\frac{5}{2}$; and since $0 = \frac{3}{4} - \frac{3}{4}$ we say
that $0 < \frac{3}{4}$.

\begin{exlist}
\item Use the first definition above to prove that $0 > -\frac{4}{3}$.
\item Use the second definition to prove that $-\frac{4}{3} < 0$.
\item Use the two definitions to prove that if $a > b$ then $b < a$.
\item If $a = b + 1$ and $b = c + 2$, prove that $a > c$.
\item Prove that if $a > b$ and $b > c$, then $a > c$.
\item If $a = b - 2$ and $b = c - 3$, then prove that $a < c$.
\item Prove that if $a < b$ and $b < c$, then $a < c$.
\item Show that if $a = b + 3$, then $3a > 3b$ and $-2a < -2b$.
\item Prove that if $a > b$, then $ca > cb$ if $c$ is positive and $ca < cb$ if
      $c$ is negative.
\item Prove that if $a < b$, then $ca < cb$ if $c$ is positive and $ca > cb$ if
      $c$ is negative.
\end{exlist}

\end{document}
